% !TEX root = ../main.tex

\section{Colaboración y Asociatividad}
\label{sec:asociatividad}

Para garantizar la viabilidad técnica, la seguridad y la correcta integración institucional del proyecto Campus Zero, la iniciativa establece una red de apoyo directa con la Dirección de Infraestructura de la Universidad. El rol estratégico de este actor clave consiste en otorgar las autorizaciones técnicas necesarias para intervenir de forma segura los circuitos secundarios de iluminación y extracción en las salas de estudio piloto, además de supervisar que el cableado y los componentes cumplan con la normativa eléctrica interna del campus. Como contraparte y beneficio mutuo, la Dirección de Infraestructura recibirá un sistema operativo funcionando, los equipos debidamente inventariados y un manual técnico estandarizado que les permitirá evaluar la posterior replicabilidad de la solución en el resto de los edificios institucionales.


\noindent\textbf{Carta de Apoyo Departamento de Infraestructura:}

\begin{quote}
    \begin{flushright}
        Valparaíso, 5 de noviembre de 2026
    \end{flushright}

    \noindent\textbf{A la Comisión Evaluadora de Proyectos Estudiantiles} \\
    \noindent\textbf{Universidad Técnica Federico Santa María}

    \vspace{0.8em}

    \noindent Por medio de la presente, la \textbf{Dirección de Infraestructura de la Universidad Técnica Federico Santa María} manifiesta su conocimiento, interés y respaldo a la iniciativa estudiantil titulada \textbf{``Campus Zero: Control Automatizado de Iluminación y Climatización de Salas''}.

    \vspace{0.8em}

    \noindent Comprendemos que el proyecto tiene como propósito principal mitigar el consumo fantasma de energía eléctrica mediante la instalación de un sistema automatizado piloto basado en sensores de presencia en diez salas de estudio del campus. 

    \vspace{0.8em}

    \noindent En el marco de nuestra colaboración institucional, facilitaremos los permisos de acceso a las dependencias involucradas, supervisaremos las directrices de seguridad operacional durante el montaje de los dispositivos y colaboraremos en la recepción formal de los equipos inventariados y de la guía técnica resultante, insumos de gran valor para evaluar futuras políticas de eficiencia energética en nuestra casa de estudios.

    \vspace{0.8em}

    \noindent Sin otro particular, saluda atentamente a ustedes,

    \vspace{1em}

    \noindent \underline{\hspace{6cm}} \\
    \noindent \textbf{Oscar Jalil - Director} \\
    \noindent Dirección de Infraestructura \\
    \noindent Universidad Técnica Federico Santa María
\end{quote}
